\documentclass[11pt,a4paper]{article}

\usepackage[T1]{fontenc}
\usepackage[utf8]{inputenc}
\usepackage[portuguese]{babel}
\usepackage{geometry}
\usepackage{setspace}
\usepackage{booktabs}
\usepackage{tabularx}
\usepackage{enumitem}
\usepackage{hyperref}
\usepackage{titlesec}

\geometry{margin=2.4cm}
\onehalfspacing
\setlist[itemize]{topsep=3pt,itemsep=2pt}
\setlist[enumerate]{topsep=3pt,itemsep=2pt}
\titleformat{\section}{\large\bfseries}{\thesection.}{0.6em}{}
\titleformat{\subsection}{\normalsize\bfseries}{\thesubsection}{0.6em}{}
\hypersetup{colorlinks=true,linkcolor=black,urlcolor=blue,citecolor=black}

\title{Estimativa de Risco de Lesões em Jogadores de Futebol\\
a partir de Dados Públicos de Desempenho:\\
uma Abordagem Baseada em Aprendizado de Máquina}
\author{Eduardo Henrique Queiroz Almeida\\
\texttt{eduardo.almeida@aluno.cefetmg.br}
\and
Sergio Henrique Quedas Ramos\\
\texttt{sergio.ramos@aluno.cefetmg.br}}
\date{Pré-projeto --- Ciência de Dados\\
CEFET-MG \quad \today}

\begin{document}
\maketitle

\begin{abstract}
\noindent
Este pré-projeto apresenta o plano de um trabalho em dupla cujo objetivo é estimar o risco de lesão de jogadores profissionais de futebol usando apenas dados públicos de partidas --- minutos jogados, congestionamento de calendário e histórico de lesões --- sem acesso a GPS ou fisiologia de clube. A fonte é o Transfermarkt (Brasileirão Série~A, temporadas 2022--2024). A metodologia combina extração via \textit{web scraping}, limpeza e engenharia de \textit{features} de carga (incluindo a razão aguda:crônica, ACWR), análise exploratória e comparação de modelos (Regressão Logística, Random Forest/XGBoost e rede neural MLP), com validação temporal e métricas adequadas a classes desbalanceadas. O resultado esperado é um \textit{score} de risco relativo, não um diagnóstico médico.
\end{abstract}

\tableofcontents
\vspace{1em}
\noindent\textit{Documento escrito para a apresentação de 10 a 15 minutos. Os slides devem espelhar as seções abaixo.}

% ---------------------------------------------------------------------------
\section{Introdução}

Lesões estão entre os fatores mais custosos do futebol profissional: afastam o atleta, prejudicam o rendimento do clube e reduzem o valor de mercado do jogador. A gestão de carga costumava ser empírica; com dados de desempenho e técnicas de Aprendizado de Máquina, passou a ser possível estimar risco de forma mais objetiva.

A literatura de ciência do esporte associa o acúmulo de carga física --- minutos, densidade de jogos, histórico prévio --- a maior probabilidade de lesão em janelas seguintes~\cite{rossi2022,applsci2024}. Clubes de elite usam GPS e sensores fisiológicos para isso. Esses dados são proprietários e inacessíveis num projeto acadêmico.

A hipótese deste trabalho, alinhada a discussões da indústria e da literatura~\cite{sportmonks}, é que dados públicos de partida já constituem um \textit{proxy} razoável de carga. O recorte proposto é o Brasileirão Série~A (2022--2024), próximo da fonte e da temporada usadas em estudo acadêmico recente (UNICESUMAR / XIV~EPCC), o que permite checagem de sanidade dos números extraídos.

O projeto percorre o ciclo completo: extração e integração dos dados, engenharia de \textit{features} de carga, análise exploratória e modelagem preditiva com validação temporal e interpretabilidade.

% ---------------------------------------------------------------------------
\section{Descrição do problema}

\subsection{Pergunta de pesquisa}

\begin{quote}
\textit{É possível estimar, de forma útil, o risco de um jogador se lesionar numa janela futura usando apenas dados públicos de partidas --- sem GPS, frequência cardíaca ou carga de treino?}
\end{quote}

\subsection{Formulação como aprendizado supervisionado}

Cada linha do dataset é uma partida da Série~A em que o jogador entrou em campo. O alvo $y$ é binário:

\begin{quote}
$y = 1$ se o jogador sofre uma lesão (não doença / gestão de esforço) nos próximos 14 dias; $y = 0$ caso contrário.
\end{quote}

As variáveis explicativas são indicadores de carga derivados do histórico de minutos e do calendário, além do perfil do atleta (idade, posição, lesões prévias).

\subsection{Por que o problema é difícil}

\begin{itemize}
  \item \textbf{Classe rara:} lesão é evento pouco frequente em relação ao número de partidas. Um modelo que sempre prevê ``sem lesão'' teria acurácia alta e pouco valor --- por isso a acurácia não será a métrica principal.
  \item \textbf{\textit{Proxy} incompleto:} minutos de jogo aproximam a carga real, mas ignoram treinos, intensidade e qualidade do preparo físico do clube.
  \item \textbf{Viés de registro:} fontes públicas tendem a ficar mais completas em anos recentes, o que pode alterar a taxa observada entre temporadas e importa no \textit{split} temporal.
  \item \textbf{Vazamento temporal:} embaralhar partidas do mesmo jogador entre treino e teste invalidaria a avaliação. O corte precisa ser no tempo.
\end{itemize}

O produto não é a frase ``o jogador X vai se lesionar'', e sim um \textit{score} de risco relativo associado a um padrão de carga.

% ---------------------------------------------------------------------------
\section{Objetivos}

\subsection{Objetivo geral}

Desenvolver e comparar modelos de Aprendizado de Máquina capazes de estimar a probabilidade de um jogador profissional sofrer uma lesão nos 14 dias seguintes a uma partida, com base em indicadores de carga derivados de dados públicos.

\subsection{Objetivos específicos}

\begin{enumerate}
  \item Construir uma base integrada (elenco, histórico de lesões e desempenho por partida) para todos os clubes da Série~A nas temporadas 2022, 2023 e 2024.
  \item Limpar os dados e categorizar o tipo de lesão (taxonomia anatômica: coxa, joelho, tornozelo, panturrilha, adutores etc.), excluindo ausências não traumáticas do alvo.
  \item Derivar \textit{features} de carga usadas na literatura: minutos agudos (7 dias), carga crônica (média de 4 semanas), ACWR (razão aguda:crônica), dias desde a última partida, jogos nos últimos 30 dias e histórico prévio de lesões.
  \item Realizar análise exploratória: perfil das lesões, calendário/congestionamento, relação carga $\times$ lesão e comparativos por clube e posição.
  \item Comparar Regressão Logística (baseline), Random Forest e/ou XGBoost e uma rede neural MLP, com tratamento de desbalanceamento e validação temporal.
  \item Interpretar os modelos (SHAP ou importância de variáveis) e discutir criticamente as limitações do \textit{proxy} público.
\end{enumerate}

% ---------------------------------------------------------------------------
\section{Metodologia}

O fluxo previsto é: extração $\rightarrow$ cache do bruto $\rightarrow$ limpeza e \textit{features} $\rightarrow$ análise exploratória $\rightarrow$ modelagem. As bibliotecas podem ser ajustadas ao longo da pesquisa.

\subsection{Técnica de extração de dados}

Não há API oficial do Transfermarkt. A extração é \textit{web scraping} com \texttt{requests} e \texttt{BeautifulSoup4} (\texttt{lxml}), em Python~3, com as seguintes práticas:

\begin{itemize}
  \item \textit{User-Agent} de navegador, atraso aleatório entre requisições (2,5--5\,s) e \textit{retry} com \textit{backoff} em HTTP 403/429;
  \item cache do HTML/JSON em \texttt{data/raw/}, para não rebaixar se o \textit{parser} quebrar;
  \item processamento em lotes de 50 jogadores, com \textit{checkpoint} (\texttt{--retomar}).
\end{itemize}

Duas fontes por jogador:
\begin{itemize}
  \item \textbf{Lesões:} aba \texttt{/verletzungen/} --- HTML estático (tabela \texttt{items}), inclusive paginação.
  \item \textbf{Desempenho:} a aba visual é renderizada em JavaScript; os dados vêm da API interna \texttt{/ceapi/performance-game/<id>}, com uma linha por partida (data, competição, minutos, titular, idade, estado de participação).
\end{itemize}

O elenco de cada clube/temporada sai da página \texttt{/startseite/verein/...}. O site usa convenção europeia de temporada: o Brasileirão de 2024 corresponde a \texttt{saison\_id=2023}. Os \textit{scripts} recebem o ano civil e convertem internamente.

\subsection{Tratamento de dados: natureza e tipos}

Os dados são observacionais, longitudinais (várias partidas por jogador) e de natureza mista:

\begin{center}
\begin{tabular}{lll}
\toprule
\textbf{Campo} & \textbf{Tipo} & \textbf{Papel} \\
\midrule
\texttt{player\_id}, \texttt{game\_id}, \texttt{clube\_id} & categórico (ID) & chave \\
\texttt{data} & temporal (\textit{datetime}) & ordem e janelas \\
\texttt{minutos\_jogados}, idade & numérico contínuo/discreto & carga e perfil \\
\texttt{participacao}, posição, local & categórico & contexto \\
\texttt{tipo\_lesao} (texto livre) & qualitativo $\rightarrow$ categoria & taxonomia \\
\texttt{y} & binário & alvo \\
\bottomrule
\end{tabular}
\end{center}

Tratamentos previstos:
\begin{itemize}
  \item tipagem de datas e IDs; minutos $=0$ se o jogador não atuou;
  \item exclusão do alvo de ausências não traumáticas (COVID, gripe, gestão de esforço, recuperação de ritmo etc.);
  \item categorização anatômica do texto livre (coxa, joelho, tornozelo, \ldots);
  \item \textit{features} de carga calculadas só com o histórico anterior à partida (sem vazar a observação atual);
  \item recorte de modelagem: \texttt{participacao=played}, competição \texttt{BRA1}, anos 2022--2024.
\end{itemize}

O ACWR é a razão entre minutos nos últimos 7 dias e a média semanal dos últimos 28 dias. Valores próximos de 1 indicam carga recente alinhada ao habitual; valores altos indicam pico recente --- associação clássica com risco na literatura~\cite{rossi2022,iran2025}.

\subsection{Visualização de dados}

A análise exploratória seguirá cinco eixos:

\begin{enumerate}
  \item panorama da base (volumes, nulos, desbalanceamento e evolução do registro ao longo dos anos);
  \item perfil das lesões (anatomia, posição, idade, duração);
  \item calendário e congestionamento (taxa ao longo da temporada, densidade de jogos);
  \item carga $\times$ lesão (ACWR, minutos agudos/crônicos, histórico prévio);
  \item comparativos entre clubes e grupos de posição.
\end{enumerate}

Bibliotecas: \texttt{pandas} (tabelas), \texttt{matplotlib} e \texttt{seaborn} (gráficos). Ferramentas de apoio: Jupyter Notebook e o repositório Git do projeto. Se necessário, outras bibliotecas de visualização poderão ser incluídas depois.

\subsection{Técnicas de machine learning}

Problema de classificação binária desbalanceada. Modelos a comparar:

\begin{itemize}
  \item \textbf{Regressão Logística} --- baseline interpretável;
  \item \textbf{Random Forest} e/ou \textbf{XGBoost} --- ensembles de árvores; a literatura em dados tabulares de lesão/carga costuma favorecê-los~\cite{applsci2024};
  \item \textbf{MLP} (rede neural, \texttt{scikit-learn} ou PyTorch) --- para avaliar se a complexidade extra se justifica em dado tabular.
\end{itemize}

Desbalanceamento: \texttt{class\_weight} e, se couber, reamostragem (SMOTE ou \textit{undersampling}).

Validação: \textit{split} temporal (treinar em 2022--2023, testar em 2024), sem embaralhar partidas do mesmo jogador. Validação cruzada agrupada por jogador pode complementar.

Métricas: AUC-ROC, AUC-PR e F1 da classe positiva. Acurácia não será a métrica principal.

Interpretabilidade: valores SHAP (ou importância nativa das árvores), para contrastar com a literatura (carga aguda, congestionamento, histórico prévio).

Implementação prevista: \texttt{scikit-learn}, \texttt{xgboost}, \texttt{imbalanced-learn} (se SMOTE) e \texttt{shap}.

% ---------------------------------------------------------------------------
\section{Repositório}

Dois sentidos de ``repositório'' importam para a pesquisa.

\subsection{Fonte dos dados (repositório externo)}

O repositório de dados é o Transfermarkt\footnote{\url{https://www.transfermarkt.com.br/}}, base pública mais usada em estudos acadêmicos de lesão quando não há alternativa aberta equivalente. Recorte:

\begin{itemize}
  \item liga: Campeonato Brasileiro Série~A (\texttt{BRA1});
  \item temporadas civis: 2022, 2023 e 2024 (20 clubes por ano, com variação natural por acesso e rebaixamento);
  \item unidade de análise no modelo: partida jogada na liga.
\end{itemize}

Não redistribuímos o HTML bruto nem os CSVs grandes. A extração é acadêmica, com \textit{rate limiting}.

\subsection{Repositório do projeto (código e análises)}

O código e os \textit{notebooks} estão em:

\begin{center}
\url{https://github.com/edualmeidahr/injury-risk-football-ml}
\end{center}

Versionam-se apenas \textit{scripts}, \textit{notebooks} e documentação. \texttt{data/raw/} e \texttt{data/processed/} entram no \texttt{.gitignore}. Estrutura principal:

\begin{itemize}
  \item \texttt{scripts/} --- cliente HTTP, \textit{scrapers} (liga, elenco, lesões, desempenho), \texttt{features.py} e \texttt{build\_dataset.py};
  \item \texttt{notebooks/} --- análise exploratória e modelagem;
  \item \texttt{requirements.txt} --- ambiente reproduzível (\texttt{pandas}, \texttt{requests}, \texttt{beautifulsoup4}, \texttt{matplotlib}, \texttt{seaborn}; depois \texttt{scikit-learn}, \texttt{xgboost} etc.).
\end{itemize}

% ---------------------------------------------------------------------------
\section{Conclusão e pretensões}

Este pré-projeto fixa o problema (risco de lesão em 14 dias a partir de carga pública), o recorte (Série~A 2022--2024) e o caminho metodológico (extração responsável, \textit{features} de carga, EDA e comparação de três famílias de modelos sob validação temporal).

Pretensões do trabalho:
\begin{itemize}
  \item entregar uma base documentada, reutilizável em trabalhos futuros;
  \item mostrar, com evidência empírica, se o \textit{proxy} público é informativo o bastante para um \textit{score} de risco relativo;
  \item comparar de forma honesta regressão, ensembles e MLP --- inclusive se a rede neural se justifica em dado tabular e desbalanceado;
  \item identificar, via interpretabilidade, quais variáveis mais se associam ao risco nesta amostra, à luz da literatura;
  \item deixar explícitas as limitações (sem GPS/treino, viés de registro, confundidores de clube) e o caráter não diagnóstico do modelo.
\end{itemize}

O trabalho pretende, ao final, um relatório que una a base construída, a exploração dos dados e a comparação dos modelos, com discussão explícita das limitações.

\begin{thebibliography}{9}

\bibitem{rossi2022}
A.~Rossi, L.~Pappalardo e P.~Cintia.
\newblock Machine Learning for Understanding and Predicting Injuries in Football.
\newblock \textit{Sports Medicine -- Open}, 2022.
\newblock \url{https://doi.org/10.1186/s40798-022-00465-4}

\bibitem{applsci2024}
\textit{Football Analytics: Assessing the Correlation between Workload, Injury and Performance of Football Players in the English Premier League}.
\newblock \textit{Applied Sciences}, 2024.
\newblock \url{https://doi.org/10.3390/app14167217}

\bibitem{iran2025}
\textit{Prediction of football injuries using GPS-based data in Iranian professional football players: a machine learning approach}.
\newblock \textit{Frontiers in Sports and Active Living}, 2025.
\newblock \url{https://www.frontiersin.org/articles/10.3389/fspor.2025.1425180}

\bibitem{medrxiv2025}
\textit{Informed Injury Prediction in Elite Football: Decision Theory meets Machine Learning}.
\newblock medRxiv, 2025 (preprint).
\newblock \url{https://www.medrxiv.org/content/10.1101/2025.04.23.25326218v1}

\bibitem{sportmonks}
Sportmonks.
\newblock \textit{Building an Injury Risk Model With Football Data}.
\newblock \url{https://www.sportmonks.com/blogs/building-an-injury-risk-model-with-football-data/}

\end{thebibliography}

\end{document}
